%% kapa: exam=1 round=34 year=2023 period=1 form=A subject=civil_law subject_ko=민법 seq=20 number=20
%% source: Q-Net 감정평가사 기출문제 (cst00309) artlSeq=5211789 file=raw/1008/5211789/2023년도 제34회 감정평가사 1차시험 1교시.pdf
%% conversion: kordoc 4.24.51 -> tools/gen1.py (auto); review=proofread-sonnet; flags=none
\begin{question}{20}
甲의 乙에 대한 채권의 소멸시효 완성을 독자적으로 원용할 수 있는 자를 모두 고른 것은? (다툼이 있으면 판례에 따름)
\begin{framedtext}
ㄱ. 甲이 乙에 대한 채권을 보전하기 위하여 행사한 채권자취소권의 상대방이 된 수익자 \par
ㄴ. 乙의 일반채권자 \par
ㄷ. 甲의 乙에 대한 채권을 담보하기 위한 유치권이 성립된 부동산의 매수인 \par
ㄹ. 甲의 乙에 대한 채권을 담보하기 위해 저당권이 설정된 경우, 그 후순위 저당권자 \par
\end{framedtext}
\begin{choices}
  \item ㄱ, ㄴ
  \item ㄱ, ㄷ
  \item ㄴ, ㄹ
  \item ㄱ, ㄷ, ㄹ
  \item ㄴ, ㄷ, ㄹ
\end{choices}
\end{question}
\answer{2}
